\chapter{Asian and Emerging Equity Markets}\label{ch:m1:asian-and-emerging-equity-markets}

In Shanghai a share that has risen ten percent since yesterday's close
cannot rise further today, and whoever bought it this morning cannot sell it
until tomorrow. In Hong Kong buyer and seller each hand a tenth of a percent
of every trade to the government. In Tokyo the smallest price step of a share
depends on which index it belongs to. In Mumbai the seller of an option pays
a tax on the premium, a trade may settle the same afternoon, and a foreign
fund may find that foreigners, collectively, already own as much of the
company as the law allows. A strategy that works in New York does not
``port'' to these markets; it has to be re-derived under a different set of
rules. This chapter is a field guide to the rules that differ, and a method
for reading any new market.

\section{What can differ}

\begin{method}[Reading a new equity market]\label{met:m1:asian-and-emerging-equity-markets:checklist}
Before any data is loaded, fill in this table for the market.
\begin{enumerate}
\item \textbf{Access:} who may hold shares directly; licences, quotas,
  ownership limits; whether a foreign investor trades the same line of stock
  as a domestic one.
\item \textbf{Unit and tick:} the lot in which orders must be sized; the
  tick table.
\item \textbf{Limits and halts:} daily \omterm{def:m1:asian-and-emerging-equity-markets:limit}{price limits}, circuit breakers, their
  reference price.
\item \textbf{Selling:} whether shares bought today can be sold today;
  whether, and how, they can be sold short.
\item \textbf{Settlement and cash:} cycle, pre-funding, currency
  convertibility.
\item \textbf{Taxes and fees:} per-trade taxes by instrument and by side;
  dividend withholding; capital gains.
\item \textbf{Sessions:} auctions, lunch breaks, after-hours fixed-price
  sessions.
\end{enumerate}
Each line changes which strategies can exist; several of them change the
statistical properties of the data itself.
\end{method}

\section{Price limits}

\begin{definition}[Price limit]\label{def:m1:asian-and-emerging-equity-markets:limit}
A daily \emph{price limit}\index{price limit} forbids orders and trades at
prices further than a fixed percentage (or amount) from a reference price,
usually the previous close. A share at its upper limit is \emph{limit-up}:
buy orders queue at the limit price and trade only if someone sells there.
\end{definition}

\begin{proposition}[Days locked at the limit]\label{prop:m1:asian-and-emerging-equity-markets:locked}
News changes the fair value of a share by a factor $1+S$ with $S>0$. Under a
daily limit $L$, the share spends
\[
n = \left\lceil \frac{\ln(1+S)}{\ln(1+L)} \right\rceil - 1
\]
full days locked at the upper limit before a day on which it can trade at its
fair value. For a fall, replace $1+S$ and $1+L$ by $1-S$ and $1-L$.
\end{proposition}
\begin{proof}
After $k$ limit days the price is at most $(1+L)^k$ times the initial one.
The fair value becomes reachable on the first day $k$ with $(1+L)^k \ge 1+S$;
the $k-1$ preceding days are locked.
\end{proof}

\begin{example}[Falls take longer]\label{ex:m1:asian-and-emerging-equity-markets:asym}
With $L = 10\%$ a rise of 50\% needs $\lceil 4.25\rceil = 5$ days, four of
them locked. A fall of 50\% needs $\lceil \ln 0.5/\ln 0.9\rceil = \lceil
6.58\rceil = 7$ days, six locked: for six days holders who want to sell
cannot, at any price. A fund that promises its own investors daily liquidity
and holds such shares has a problem that no risk model built on daily
returns will have measured (\cref{fig:m1:asian-and-emerging-equity-markets:locked}).
\end{example}

\begin{omfigure}
\begin{tikzpicture}
\begin{axis}[width=10.5cm,height=4.6cm,
  xlabel={size of the upward revaluation (\%)},ylabel={days locked limit-up},
  tick label style={font=\small},label style={font=\small},
  xmin=0,xmax=80,ymin=0,ymax=12.5,grid=major,grid style={gray!25},
  legend columns=3,legend style={font=\footnotesize,at={(0.5,-0.36)},anchor=north,draw=none,fill=none,/tikz/every even column/.append style={column sep=8pt}},
  table/col sep=comma]
\addplot[omThm,very thick,const plot] table[x=shock_pct,y=l5]{figdata/markets-1/12-asian-and-emerging-equity-markets/locked.csv};
\addplot[omDef,very thick,const plot] table[x=shock_pct,y=l10]{figdata/markets-1/12-asian-and-emerging-equity-markets/locked.csv};
\addplot[omMeth,very thick,const plot] table[x=shock_pct,y=l20]{figdata/markets-1/12-asian-and-emerging-equity-markets/locked.csv};
\legend{limit 5\%,limit 10\%,limit 20\%}
\end{axis}
\end{tikzpicture}

\omcaption{Full days spent at the limit before a revaluation can be traded,
by \cref{prop:m1:asian-and-emerging-equity-markets:locked}. Data: computed by
the chapter's script.}\label{fig:m1:asian-and-emerging-equity-markets:locked}
\end{omfigure}

\begin{proposition}[What limits do to the data]\label{prop:m1:asian-and-emerging-equity-markets:data}
Let true daily returns be independent. Returns observed under a \omterm{def:m1:asian-and-emerging-equity-markets:limit}{price limit}
(i) have a smaller variance than the true ones, (ii) are positively
autocorrelated, and (iii) have point masses at $\pm L$.
\end{proposition}
\begin{proof}[Partial proof]
(iii) is the definition. For (i) and (ii): on a limit day the observed return
is $L$ and the unexpressed part of the move is carried to the next day, where
it adds to that day's independent return. The carried part has the sign of
today's return, so consecutive observed returns covary positively; and since
a large move is split over several days, the sum of squared daily returns is
smaller than the square of the move. A full computation needs the return
distribution and is left to the tutorial's simulation.
\end{proof}

\begin{omfigure}
\begin{tikzpicture}
\begin{axis}[width=11cm,height=5cm,ybar=0pt,bar width=3.4pt,ymode=log,log origin=infty,
  xlabel={daily return (\%)},ylabel={days out of 60\,000 (log scale)},
  tick label style={font=\small},label style={font=\small},
  xmin=-14,xmax=14,ymin=1,ymax=30000,
  legend columns=2,legend style={font=\footnotesize,at={(0.5,-0.32)},anchor=north,draw=none,fill=none,/tikz/every even column/.append style={column sep=8pt}},
  area legend,table/col sep=comma]
\addplot[fill=omExo!50,draw=omExo] table[x=centre_pct,y=true]{figdata/markets-1/12-asian-and-emerging-equity-markets/limit_hist.csv};
\addplot[fill=omDef!70,draw=omDef] table[x=centre_pct,y=observed]{figdata/markets-1/12-asian-and-emerging-equity-markets/limit_hist.csv};
\legend{true returns,observed under a 10\% limit}
\end{axis}
\end{tikzpicture}

\omcaption{Heavy-tailed returns with and without a 10\% daily limit. The
observed series has no tail beyond the limit and a spike at it; its standard
deviation is 3.03\% against 3.45\%, and its first autocorrelation $+0.09$
against zero. Data: the tutorial's simulation.}\label{fig:m1:asian-and-emerging-equity-markets:hist}
\end{omfigure}

A volatility estimated from limit-constrained closes is too low; a momentum
signal fitted on them finds a ``continuation'' that cannot be traded, because
on the days it predicts nobody can buy. Both are among the commonest errors
in cross-market research.

\section{Mainland China and the Connect}

\begin{definition}[T+0 restriction]\label{def:m1:asian-and-emerging-equity-markets:tplus0}
A \emph{T+0 restriction}\index{T+0 restriction} (in mainland China, the
``T+1 trading rule'') forbids selling shares on the day they were bought. It
is a rule about trading, distinct from the \omterm{def:m1:clearing-and-settlement:cycle}{settlement cycle}.
\end{definition}

\begin{definition}[Qualified foreign investor and Stock Connect]\label{def:m1:asian-and-emerging-equity-markets:connect}
A \emph{qualified foreign investor}\index{qualified foreign investor} is a
foreign institution licensed by the mainland regulator to invest directly in
onshore securities. \emph{Stock Connect}\index{Stock Connect} is the link
between the Hong Kong exchange and the Shanghai and Shenzhen exchanges that
lets investors on each side trade eligible shares on the other through their
home brokers and clearing houses, within a daily quota, without a licence.
\end{definition}

\begin{dated}{2026-09}{Mainland A-share trading rules}\label{dat:m1:asian-and-emerging-equity-markets:china}
Main-board shares in Shanghai and Shenzhen may move 10\% a day; shares on the
STAR Market and ChiNext 20\%, with no limit during a STAR listing's first five
days. Revised trading rules of the Shanghai, Shenzhen and Beijing exchanges
took effect on 6 July 2026: the limit for main-board shares under
special treatment was widened to 10\%, and the after-hours fixed-price session
(15:05 to 15:30) was extended to all A-shares and exchange-traded funds.
Shares bought cannot be sold the same day. The northbound \omterm{def:m1:asian-and-emerging-equity-markets:connect}{Stock Connect} quota
is RMB 52 billion a day for each of the Shanghai and Shenzhen links.
\end{dated}

\begin{omfigure}
\begin{tikzpicture}[font=\small,
  box/.style={draw=omDef,rounded corners=2pt,minimum width=3.0cm,minimum height=1.0cm,align=center,fill=omDef!4},
  mid/.style={draw=omThm,rounded corners=2pt,minimum width=3.0cm,minimum height=1.0cm,align=center,fill=omThm!5},
  flow/.style={-{Stealth[length=2mm]},thick}]
\node[box] (f) at (0,0) {International\\investor};
\node[box] (hb) at (4.2,0) {Hong Kong\\broker};
\node[mid] (hk) at (8.4,0) {Hong Kong exchange\\and clearing house};
\node[box] (sh) at (12.8,0.9) {Shanghai\\exchange};
\node[box] (sz) at (12.8,-0.9) {Shenzhen\\exchange};
\draw[flow,omDef] (f) -- (hb);
\draw[flow,omDef] (hb) -- (hk);
\draw[flow,omDef] (hk) -- (sh);
\draw[flow,omDef] (hk) -- (sz);
\node[font=\footnotesize,omThm,align=center] at (8.4,-1.25) {daily quota checked here;\\settles in renminbi};
\end{tikzpicture}

\omcaption{A northbound \omterm{def:m1:asian-and-emerging-equity-markets:connect}{Stock Connect} order. The investor faces only its
Hong Kong broker; the Hong Kong exchange's subsidiary is the member of the
mainland exchange, and the two clearing houses are each other's
counterparties. Mainland rules --- limits, lots, the ban on same-day selling
--- apply to the order in full.}
\end{omfigure}

\begin{remark}[The same company at two prices]\label{rem:m1:asian-and-emerging-equity-markets:ah}
Many mainland companies are listed both as A-shares onshore and as H-shares
in Hong Kong. The two lines carry the same dividends and votes and are
\emph{not} convertible into each other; they routinely trade at different
prices, because the populations allowed to hold each are different and no
arbitrageur can deliver one against the other. A price difference that cannot
be closed by delivery is a fact about access, not a free lunch
(\cref{ch:m1:delta-one-instruments}).
\end{remark}

\section{Hong Kong}

\begin{definition}[Board lot]\label{def:m1:asian-and-emerging-equity-markets:lot}
A \emph{board lot}\index{board lot} is the number of shares in which orders
on an exchange's main order book must be multiples. Where the issuer chooses
it, the board lot differs from stock to stock.
\end{definition}

\begin{definition}[Stamp duty]\label{def:m1:asian-and-emerging-equity-markets:stamp}
A \emph{stamp duty}\index{stamp duty} on share transfers is a tax charged as a
percentage of the value of each transaction, on the buyer, the seller or
both, whatever the profit or loss.
\end{definition}

\begin{dated}{2026-09}{Hong Kong: duty and lots}\label{dat:m1:asian-and-emerging-equity-markets:hk}
\omterm{def:m1:asian-and-emerging-equity-markets:stamp}{Stamp duty} on Hong Kong stock transfers is 0.1\% of the value for the buyer
and 0.1\% for the seller, since 17 November 2023 (previously 0.13\%). \omterm{def:m1:asian-and-emerging-equity-markets:lot}{Board
lots} were chosen freely by issuers, in more than forty different sizes; a
reform implemented from July 2026 restricts them to eight standard sizes (1,
50, 100, 500, 1\,000, 2\,000, 5\,000 and 10\,000 shares), with a lower minimum
and a new maximum value per lot.
\end{dated}

\begin{example}[Twenty basis points]\label{ex:m1:asian-and-emerging-equity-markets:duty}
A round trip in a Hong Kong share costs 20 basis points in duty before any
spread, commission or impact. A market-making strategy that captures one or
two basis points per round trip
(\cref{dat:m1:what-a-trading-firm-does:listed}) cannot exist under that tax;
strategies with holding periods of days, earning fifty basis points per
position, can. Transaction taxes do not shave every strategy equally: they
delete the short-horizon end of the spectrum, and with it the participants
who would have provided intraday liquidity.
\end{example}

\section{Japan, Korea and the uptick rule}

The Tokyo exchange trades shares in units of 100 and publishes several tick
tables: the constituents of its large-capitalisation indices use a finer
table than other shares, down to a fraction of a yen for the most liquid. A
share that enters or leaves the index changes table, which moves its spread,
its queue lengths and the economics of making a market in it, with no change
in the company.

\begin{definition}[Uptick rule]\label{def:m1:asian-and-emerging-equity-markets:uptick}
An \emph{uptick rule}\index{uptick rule} restricts the prices at which a \omterm{def:m1:financing:short}{short
sale} may be entered: typically not below the last traded price, or not below
it once the share has fallen by a given amount during the day.
\end{definition}

\begin{example}[Seventeen months without short selling]\label{ex:m1:asian-and-emerging-equity-markets:korea}
In November 2023 the Korean authorities banned short selling of all listed
shares, citing naked \omterm{def:m1:financing:short}{short sales} by global banks. The ban ended on 31 March
2025, when short selling resumed for every listed share together with a new
electronic system for detecting naked \omterm{def:m1:financing:short}{short sales}. For seventeen months no
long--short equity strategy could be run onshore in one of the world's
larger markets; index arbitrage, convertible-bond hedging and market making
in single-stock derivatives were all impaired at once. Regulatory risk of
this kind is a parameter of the market, like volatility.
\end{example}

\section{India}

\begin{definition}[Securities transaction tax and foreign ownership limit]\label{def:m1:asian-and-emerging-equity-markets:stt}
A \emph{securities transaction tax}\index{securities transaction tax} (STT)
is a tax on exchange transactions whose rate and payer depend on the
instrument: shares, futures or options. A
\emph{foreign ownership limit}\index{foreign ownership limit} caps the
fraction of a company that non-residents may hold in aggregate; when it is
reached, foreigners can buy only from other foreigners.
\end{definition}

\begin{dated}{2026-09}{India: tax and settlement}\label{dat:m1:asian-and-emerging-equity-markets:india}
STT on shares bought or sold for delivery is 0.1\% of the value on each side.
From 1 April 2026 the rate on futures is 0.05\% of the value, paid by the
seller (previously 0.02\%), and on options 0.15\% of the premium, paid by the
seller (previously 0.10\%). Shares settle at $T{+}1$; an optional same-day
($T{+}0$) cycle is available in the 500 largest shares.
\end{dated}

\begin{omfigure}
\begin{tikzpicture}
\begin{axis}[width=9cm,height=3.6cm,xbar,bar width=8pt,
  xlabel={transaction tax on a round trip (basis points of value)},
  ytick=data,yticklabels from table={figdata/markets-1/12-asian-and-emerging-equity-markets/taxes.csv}{label},
  yticklabel style={font=\small},tick label style={font=\small},label style={font=\small},
  y dir=reverse,xmin=0,xmax=25,enlarge y limits=0.25,grid=major,grid style={gray!25},
  nodes near coords,nodes near coords style={font=\footnotesize},
  table/col sep=comma]
\addplot[fill=omDef!60,draw=omDef] table[x=bp,y=k]{figdata/markets-1/12-asian-and-emerging-equity-markets/taxes.csv};
\end{axis}
\end{tikzpicture}

\omcaption{Per-trade taxes on a purchase followed by a sale. A market with no
such tax would show zero. Data: the rates of
\cref{dat:m1:asian-and-emerging-equity-markets:hk,dat:m1:asian-and-emerging-equity-markets:india}.}
\end{omfigure}

\section{What changes for a trading firm}

\begin{proposition}[The tax floor on turnover]\label{prop:m1:asian-and-emerging-equity-markets:floor}
A strategy earns a gross edge of $e$ per round trip and pays a round-trip
transaction tax $\tau$ and other costs $\kappa$, all as fractions of value.
It is viable only if $e > \tau + \kappa$; and if its edge per unit of holding
time is roughly constant, $e \approx \alpha h$ for a holding period $h$, the
shortest viable holding period is $h^\star = (\tau+\kappa)/\alpha$.
\end{proposition}
\begin{proof}
Net profit per round trip is $e - \tau - \kappa$; substitute $e = \alpha h$.
\end{proof}

With $\alpha$ of 5 basis points a day and $\kappa$ of 5 basis points, a
market without tax admits holding periods from one day; a market with 20
basis points of duty only from five days. The same signal that is traded
intraday in one market must be slowed down, or dropped, in another.

\section{Tutorial: what a price limit does to returns}\label{tut:m1:asian-and-emerging-equity-markets:limits}

\begin{tutorial}
\textbf{Goal.} Pass a series of true returns through a 10\% limit and measure
what it does to volatility and autocorrelation. \textbf{End state:}
\cref{fig:m1:asian-and-emerging-equity-markets:hist} and its three
statistics.
\begin{tutsteps}
\item \textbf{The limit with carry-over.} The observed price cannot move more
than $L$ a day; the rest of the move waits.
\omcode{code/markets-1/12-asian-and-emerging-equity-markets/python/asia_rules.py}{17}{27}{Truncation with carry-over, in log space so that the total move is conserved.}\label{lst:m1:asian-and-emerging-equity-markets:truncate}
\item \textbf{Check conservation.} A single $+35\%$ revaluation becomes
$+10\%$, $+10\%$, $+10\%$ and a remainder; the product of the observed gross
returns is 1.35.
\item \textbf{Measure.} With heavy-tailed returns of 3.45\% daily standard
deviation, 2\% of days end at a limit; the observed standard deviation is
3.03\% and the first autocorrelation $+0.088$.
\item \textbf{Days locked.}
\omcode{code/markets-1/12-asian-and-emerging-equity-markets/python/asia_rules.py}{7}{14}{The proposition as a function, with the asymmetry between rises and falls.}\label{lst:m1:asian-and-emerging-equity-markets:locked}
\end{tutsteps}
\textbf{What to change next.} Fit a one-day momentum rule to the observed
series and compute its ``profit''; then forbid trading on limit days and
recompute. Then widen the limit to 20\% and see how much of the distortion
remains.
\end{tutorial}

\section{Build: the market-rules table}\label{bld:m1:asian-and-emerging-equity-markets:rules}

\begin{build}
\bfield{Purpose} Every order the miniature firm creates is validated, before
it leaves, against the rules of the market it goes to; every backtest asks
the same table what was allowed on that date.
\bfield{Interface} \texttt{MarketRules(market, valid\_from, lot, tick\_table,
limit\_up, limit\_down, same\_day\_sell, short\_selling, settlement\_days,
tax\_buy, tax\_sell)}; \texttt{RuleBook.load(csv)};
\texttt{rules(market, date)}; \texttt{validate(order, reference\_price,
bought\_today)} returning a list of violations; \texttt{tax(side, value)}.
\bfield{Rules} Rules are versioned by date: a query for a past date returns
the rules then in force, never today's. Quantities must be multiples of the
lot; prices multiples of the tick applicable at that price and inside the
limit band around the reference price; a sale exceeding the position held at
the start of the day is rejected where same-day selling is forbidden.
\bfield{Acceptance tests} \texttt{code/firm/marketrules/tests/}: lot and tick
rounding; limit bands; the same-day-sell rejection; a rate change on a given
date applied only from that date.
\bfield{Stretch} Load the rule history of one real market from its exchange's
circulars and reproduce the dated boxes of this chapter as queries.
\end{build}

\begin{omsources}
\item KPMG, \emph{Hong Kong: Reduced stamp duty rate on stock transfers
effective 17 November 2023}; HKEX, consultation and conclusions on the board
lot framework, 2025--2026.
\item HKEX, \emph{Stock Connect: Information Booklet and FAQ}, 2024.
\item Shanghai, Shenzhen and Beijing stock exchanges, revised trading rules
effective 6 July 2026 (as reported in the financial press).
\item Financial Services Commission of Korea, \emph{Stock Short Selling to be
Fully Reinstated from March 31}, press release, 2025.
\item ClearTax, \emph{Securities Transaction Tax: latest updates}, and ICICI
Direct, \emph{STT changes in Budget 2026} (rates from 1 April 2026);
Securities and Exchange Board of India, circulars on the optional $T{+}0$
settlement cycle.
\item Japan Exchange Group, \emph{Trading Rules of Domestic Stocks: Tick
Size}.
\item K.~Kim and S.~Rhee, ``Price limit performance: evidence from the Tokyo
Stock Exchange'', \emph{Journal of Finance} 52 (1997).
\end{omsources}

\section{Exercises}

\begin{exercise}[$\star$]\label{exo:m1:asian-and-emerging-equity-markets:1}
A share closed at 24.00 with a 10\% limit and a tick of 0.01. Give today's
upper and lower limit prices. A broker's system computes the upper limit as
26.40 exactly; under what rounding convention could the exchange's be
different?
\end{exercise}

\begin{exercise}[$\star$]\label{exo:m1:asian-and-emerging-equity-markets:2}
An investor wants HKD 250\,000 of a share trading at HKD 61.30 with a \omterm{def:m1:asian-and-emerging-equity-markets:lot}{board
lot} of 500. How many shares can it buy, for what value, and what \omterm{def:m1:asian-and-emerging-equity-markets:stamp}{stamp duty}
does it pay on the purchase?
\end{exercise}

\begin{exercise}[$\star$]\label{exo:m1:asian-and-emerging-equity-markets:3}
Using \cref{dat:m1:asian-and-emerging-equity-markets:india}, compute the tax
on (a) buying and later selling INR 5 million of shares for delivery; (b)
selling index futures with a contract value of INR 20 million; (c) selling
options for a premium of INR 400\,000.
\end{exercise}

\begin{exercise}[$\star\star$]\label{exo:m1:asian-and-emerging-equity-markets:4}
Bad news cuts a share's fair value by 35\%. How many days is it locked
limit-down under a 10\% limit? Under 20\%? Under 5\%?
\end{exercise}

\begin{exercise}[$\star\star$]\label{exo:m1:asian-and-emerging-equity-markets:5}
A signal earns 4 basis points of gross edge per day held; costs other than
tax are 6 basis points per round trip. Give the shortest viable holding
period with no transaction tax, with a 20-basis-point round-trip duty, and
the net profit per round trip at a holding period of ten days in each case.
\end{exercise}

\begin{exercise}[$\star\star$]\label{exo:m1:asian-and-emerging-equity-markets:6}
A trader in a market with a ban on same-day selling buys 10\,000 shares at
the open; at noon the share is up 6\% and she wants to lock in the gain. The
share has a liquid futures contract and no single-stock options. What can she
do, and what risk remains?
\end{exercise}

\begin{exercise}[$\star\star\star$]\label{exo:m1:asian-and-emerging-equity-markets:7}
\emph{Coding.} With \texttt{truncate}, generate 60\,000 heavy-tailed returns
as in the tutorial (seed 12) and compute the observed standard deviation and
first autocorrelation for limits of 5\%, 10\% and 20\%. Explain the pattern.
\end{exercise}

\begin{exercise}[$\star\star\star$]\label{exo:m1:asian-and-emerging-equity-markets:8}
\emph{Find the flaw.} A researcher reports: ``A-shares that close \omterm{def:m1:asian-and-emerging-equity-markets:limit}{limit-up}
gain a further 3.1\% on average the next day; buying at the \omterm{def:m1:asian-and-emerging-equity-markets:limit}{limit-up} close
earns 3.1\% a day.'' Identify the two reasons the profit is not available,
and the statistic that should have been reported.
\end{exercise}

\section{Problem: Limit-Up}

\begin{problem}[{Weekend problem --- a takeover in a market with limits}]\label{pb:m1:asian-and-emerging-equity-markets:1}
A share closes at 20.00 on Monday in a market with a 10\% daily \omterm{def:m1:asian-and-emerging-equity-markets:limit}{price limit}
(limit prices are rounded down to the cent), a \omterm{def:m1:asian-and-emerging-equity-markets:lot}{board lot} of 100 shares, a ban
on selling shares bought the same day, and a duty of 0.1\% on each side. After
the close an acquirer announces a cash offer at 29.00, which the market
considers certain to complete in three months. The interest rate is 4\% a
year.

\medskip\noindent\textbf{Part I --- The path.}
\begin{enumerate}
\item What is the share worth on Tuesday morning, discounting the offer at
the interest rate for three months?
\item Give the upper limit price on Tuesday, Wednesday, Thursday and Friday,
each computed from the previous close.
\item On which day can the share first trade at its fair value? How many days
is it locked?
\item Check against \cref{prop:m1:asian-and-emerging-equity-markets:locked}.
\end{enumerate}

\medskip\noindent\textbf{Part II --- The queue.}
On Tuesday buy orders for 40 million shares queue at the limit price; 300\,000
shares are sold there during the day, allocated in time priority. You entered
an order for 50\,000 shares 0.2 seconds after the open, behind 9 million
shares already queued.
\begin{enumerate}[resume]
\item How many shares do you receive?
\item What is a share bought at Tuesday's limit worth relative to its price,
in percent?
\item Someone at the very front of the queue buys 100\,000 shares. What is
the expected gain in currency, before costs?
\item What determines who is at the front, and what would a firm pay for
that?
\item If you do buy on Tuesday, when can you first sell, and at what limit
price at best?
\end{enumerate}

\medskip\noindent\textbf{Part III --- The data.}
\begin{enumerate}[resume]
\item Write the series of daily close-to-close returns from Monday to the
first day at fair value, assuming the share then closes at fair value.
\item Give the sum of squared daily returns, and the square of the total
return; compare.
\item A risk system estimates this share's volatility from the last 20 daily
returns, the other 16 being zero. What daily volatility does it report, and
what would it report had the move occurred in one day?
\item Compute the first autocorrelation of your series of question 10
(without subtracting the mean). What would a momentum signal conclude?
\end{enumerate}

\medskip\noindent\textbf{Part IV --- The trade that exists.}
\begin{enumerate}[resume]
\item After the share reaches fair value it trades at 28.60 with the offer
at 29.00 three months away. Give the gross return of buying and tendering,
annualised.
\item Subtract the duty on the purchase (tendering is a sale) and the
financing at 4\%. Is the trade worth doing?
\item Your \omterm{def:m1:asian-and-emerging-equity-markets:lot}{board lot} is 100 shares and you want to invest 2 million. How many
shares?
\item The same company has a line listed in another market with no \omterm{def:m1:asian-and-emerging-equity-markets:limit}{price
limit}, not convertible into this one. What would you expect that line to have
done on Tuesday, and what does the gap between the two lines measure?
\item Name the one risk that the certain-completion assumption removed, and
say how it would change question 14.
\item State the \emph{named result}: the number of days the share is locked
at the limit.
\item In one sentence, what does a \omterm{def:m1:asian-and-emerging-equity-markets:limit}{price limit} protect, and what does it
cost?
\end{enumerate}
\end{problem}

\section{Interview questions}

\begin{interviewq}[$\star$ \iqroles{trader, researcher}]\label{iq:m1:asian-and-emerging-equity-markets:1}
What is a daily \omterm{def:m1:asian-and-emerging-equity-markets:limit}{price limit}, and what happens to a share when its fair value
moves by more than the limit?
\end{interviewq}

\begin{interviewq}[$\star$ \iqroles{researcher, mle}]\label{iq:m1:asian-and-emerging-equity-markets:2}
You are given daily closes for a market with 10\% \omterm{def:m1:asian-and-emerging-equity-markets:limit}{price limits}. Name two
statistics that will be biased if you treat them like US data.
\end{interviewq}

\begin{interviewq}[$\star\star$ \iqroles{trader, researcher}]\label{iq:m1:asian-and-emerging-equity-markets:3}
A market charges 10 basis points of tax on each side of a share trade. Which
strategies disappear, and who ends up providing liquidity?
\end{interviewq}

\begin{interviewq}[$\star\star$ \iqroles{researcher, trader}]\label{iq:m1:asian-and-emerging-equity-markets:4}
A company's onshore and offshore shares trade 30\% apart. Why is that not an
arbitrage, and how might you still trade it?
\end{interviewq}

\begin{interviewq}[$\star\star$ \iqroles{developer}]\label{iq:m1:asian-and-emerging-equity-markets:5}
Your order-validation layer must support fifteen markets whose lot, tick and
limit rules change several times a year. How do you design it?
\end{interviewq}

\begin{interviewq}[$\star\star\star$ \iqroles{researcher, trader}]\label{iq:m1:asian-and-emerging-equity-markets:6}
A regulator bans short selling overnight. List the strategies and market
functions affected, in order of severity, and what you would do on the first
morning if you ran a long--short book there.
\end{interviewq}
